\documentclass[12pt]{ximera}
\title{Homework 8: Apportionment --- Hamilton, Jefferson, Adams, Webster}
\begin{document}
\begin{abstract}
Sections 4.1--4.4: standard divisors and quotas, Hamilton's method applied to
apportioning RAs across seventeen residence halls, and the three divisor methods
applied to a district where one of them breaks the quota rule.
\end{abstract}
\maketitle

\section*{Sections 4.1--4.4: Methods of Apportionment}

\begin{problem}
The residence halls at a university have the following population capacities,
totaling $6558$ students:

\begin{center}
\begin{tabular}{r|c|c|c|c|c|c}
\textbf{Hall} & Aspen & Cedar & Pine & Birch & Willow & Larch \\\hline
\textbf{Pop.} & 270 & 270 & 331 & 184 & 544 & 544 \\
\end{tabular}

\smallskip
\begin{tabular}{r|c|c|c|c|c|c}
\textbf{Hall} & Spruce & Poplar & Hickory & Maple & Linden & Elm \\\hline
\textbf{Pop.} & 99 & 136 & 107 & 496 & 321 & 202 \\
\end{tabular}

\smallskip
\begin{tabular}{r|c|c|c|c|c}
\textbf{Hall} & Walnut & Oak & Alder & Chestnut & Sycamore \\\hline
\textbf{Pop.} & 770 & 235 & 116 & 713 & 1220 \\
\end{tabular}
\end{center}

The university has hired $220$ RAs altogether and wants to apportion them using
Hamilton's method.

\begin{prompt}
\textbf{(a)} Compute the standard divisor, to two decimal places.

$\answer[tolerance=0.005]{29.81}$

\smallskip
What does that value represent in this situation?
\begin{multipleChoice}
\choice[correct]{The number of students per RA}
\choice{The number of RAs per residence hall}
\choice{The number of residence halls per RA}
\end{multipleChoice}

\begin{feedback}
The standard divisor is the total population divided by the number of seats:
\[ SD = \frac{6558}{220} \approx 29.81. \]
It represents the number of \emph{students per RA} --- roughly one RA for every $30$
students.
\end{feedback}
\end{prompt}

\begin{prompt}
\textbf{(b)} Compute the standard quota of each of these residence halls, to two
decimal places. (A hall's quota is its population divided by the standard divisor.)

Sycamore: $\answer[tolerance=0.005]{40.93}$ \qquad
Walnut: $\answer[tolerance=0.005]{25.83}$ \qquad
Chestnut: $\answer[tolerance=0.005]{23.92}$

\smallskip
Willow: $\answer[tolerance=0.005]{18.25}$ \qquad
Maple: $\answer[tolerance=0.005]{16.64}$ \qquad
Spruce: $\answer[tolerance=0.005]{3.32}$

\begin{feedback}
Dividing each population by $29.81$:
\[
\begin{array}{llll}
\text{Aspen } 9.06 & \text{Cedar } 9.06 & \text{Pine } 11.10 & \text{Birch } 6.17\\
\text{Willow } 18.25 & \text{Larch } 18.25 & \text{Spruce } 3.32 & \text{Poplar } 4.56\\
\text{Hickory } 3.59 & \text{Maple } 16.64 & \text{Linden } 10.77 & \text{Elm } 6.78\\
\text{Walnut } 25.83 & \text{Oak } 7.88 & \text{Alder } 3.89 & \text{Chestnut } 23.92\\
\text{Sycamore } 40.93 & & &
\end{array}
\]
These sum to $220$, as the standard quotas always must.
\end{feedback}
\end{prompt}

\begin{prompt}
\textbf{(c)} What is the total of the (standard) \emph{lower} quotas?

$\answer{211}$

\smallskip
How many surplus RAs remain to be handed out? $\answer{9}$

\begin{hint}
Round every quota \emph{down} and add. The shortfall from $220$ is the surplus.
\end{hint}

\begin{feedback}
Rounding every quota down and adding gives $211$. Since $220$ RAs must be placed, there
is a surplus of
\[ 220 - 211 = 9 \]
RAs still to assign.
\end{feedback}
\end{prompt}

\begin{prompt}
\textbf{(d)} Under Hamilton's method, the surplus goes to the halls with the largest
fractional parts. Which halls receive an extra RA?

\begin{selectAll}
\choice{Aspen}
\choice{Cedar}
\choice{Pine}
\choice{Birch}
\choice{Willow}
\choice{Larch}
\choice{Spruce}
\choice{Poplar}
\choice[correct]{Hickory}
\choice[correct]{Maple}
\choice[correct]{Linden}
\choice[correct]{Elm}
\choice[correct]{Walnut}
\choice[correct]{Oak}
\choice[correct]{Alder}
\choice[correct]{Chestnut}
\choice[correct]{Sycamore}
\end{selectAll}

\begin{feedback}
Ranking the halls by fractional part, largest first:
\[
\begin{array}{lll}
\text{Sycamore } .927 & \text{Chestnut } .919 & \text{Alder } .891\\
\text{Oak } .884 & \text{Walnut } .831 & \text{Elm } .776\\
\text{Linden } .769 & \text{Maple } .639 & \text{Hickory } .590\\
\end{array}
\]
Those are the top nine, so each of them gets one extra RA. The next in line is Poplar
at $.562$, which just misses.

The final Hamilton apportionment is:
\[
\begin{array}{llll}
\text{Aspen } 9 & \text{Cedar } 9 & \text{Pine } 11 & \text{Birch } 6\\
\text{Willow } 18 & \text{Larch } 18 & \text{Spruce } 3 & \text{Poplar } 4\\
\text{Hickory } 4 & \text{Maple } 17 & \text{Linden } 11 & \text{Elm } 7\\
\text{Walnut } 26 & \text{Oak } 8 & \text{Alder } 4 & \text{Chestnut } 24\\
\text{Sycamore } 41 & & &
\end{array}
\]
which totals $220$. \checkmark

Notice Hickory (population $107$) receives $4$ RAs while Poplar (population $136$)
receives only $4$ as well --- and Alder, with just $116$ students, also gets $4$. Small
halls can do disproportionately well or badly on the strength of a fractional part.
\end{feedback}
\end{prompt}
\end{problem}

\begin{problem}
The city of Gotham needs to apportion $100$ new villains to its four districts:

\begin{center}
\begin{tabular}{|r|c|c|c|c|c|}\hline
\textbf{District} & A & B & C & D & \textbf{Total} \\\hline
\textbf{Population} & 86{,}915 & 5{,}400 & 4{,}325 & 3{,}360 & 100{,}000 \\\hline
\end{tabular}
\end{center}

\begin{prompt}
\textbf{(a)} Compute the standard divisor and the standard quotas.

Standard divisor: $\answer{1000}$

\smallskip
$q_A = \answer[tolerance=0.0005]{86.915}$ \quad
$q_B = \answer[tolerance=0.0005]{5.4}$ \quad
$q_C = \answer[tolerance=0.0005]{4.325}$ \quad
$q_D = \answer[tolerance=0.0005]{3.36}$

\begin{feedback}
\[ SD = \frac{100{,}000}{100} = 1000, \]
so the quotas are just the populations divided by $1000$:
\[ q_A = 86.915, \quad q_B = 5.400, \quad q_C = 4.325, \quad q_D = 3.360, \]
totaling $100$. \checkmark
\end{feedback}
\end{prompt}

\begin{prompt}
\textbf{(b)} Apply Hamilton's method.

$A = \answer{87}$ \quad $B = \answer{6}$ \quad $C = \answer{4}$ \quad $D = \answer{3}$

\begin{feedback}
The lower quotas are $86, 5, 4, 3$, totaling $98$, so $2$ surplus seats remain.
The fractional parts are
\[ A: .915, \qquad B: .400, \qquad C: .325, \qquad D: .360, \]
so the two largest --- A and B --- each gain a seat:
\[ A = 87, \quad B = 6, \quad C = 4, \quad D = 3. \]
\end{feedback}
\end{prompt}

\begin{prompt}
\textbf{(c)} Apply Jefferson's method (round every modified quota \emph{down}), using a
modified divisor between $900$ and $1100$.

$A = \answer{88}$ \quad $B = \answer{5}$ \quad $C = \answer{4}$ \quad $D = \answer{3}$

\begin{hint}
Rounding down at the standard divisor gives only $98$ seats, so you need a
\emph{smaller} divisor to push the quotas up. Try something around $980$.
\end{hint}

\begin{feedback}
Rounding down at $SD = 1000$ yields only $98$ seats, so the divisor must shrink.
Any modified divisor from about $976.6$ to $987.7$ works. Taking $d = 980$:
\[
\begin{array}{llll}
A: \frac{86915}{980}=88.69 \to 88 &
B: \frac{5400}{980}=5.51 \to 5 &
C: \frac{4325}{980}=4.41 \to 4 &
D: \frac{3360}{980}=3.43 \to 3
\end{array}
\]
which totals $100$. \checkmark

\textbf{Note the quota rule violation.} District A's standard quota was $86.915$, so
its lower and upper quotas are $86$ and $87$. Jefferson's method awards A \textbf{88}
seats --- more than its upper quota. Jefferson's method systematically favors large
districts in exactly this way.
\end{feedback}
\end{prompt}

\begin{prompt}
\textbf{(d)} Apply Adams' method (round every modified quota \emph{up}).

$A = \answer{85}$ \quad $B = \answer{6}$ \quad $C = \answer{5}$ \quad $D = \answer{4}$

\begin{hint}
Rounding up at the standard divisor gives too many seats, so you need a \emph{larger}
divisor. Try something around $1030$.
\end{hint}

\begin{feedback}
Rounding up at $SD=1000$ gives $87+6+5+4 = 102$, which is too many, so the divisor must
grow. Any modified divisor from about $1022.5$ to $1034.7$ works. Taking $d = 1030$:
\[
\begin{array}{llll}
A: \frac{86915}{1030}=84.38 \to 85 &
B: \frac{5400}{1030}=5.24 \to 6 &
C: \frac{4325}{1030}=4.20 \to 5 &
D: \frac{3360}{1030}=3.26 \to 4
\end{array}
\]
totaling $100$. \checkmark

\textbf{Another quota rule violation, in the opposite direction.} A's lower quota is
$86$, but Adams' method awards it only $85$. Adams' method favors small districts, just
as Jefferson's favors large ones.
\end{feedback}
\end{prompt}

\begin{prompt}
\textbf{(e)} Apply Webster's method (conventional rounding of the modified quotas).

$A = \answer{88}$ \quad $B = \answer{5}$ \quad $C = \answer{4}$ \quad $D = \answer{3}$

\begin{feedback}
Conventional rounding at $SD = 1000$ gives $87+5+4+3 = 99$, one short, so the divisor
must shrink slightly. Any modified divisor from about $982.1$ to $993.3$ works. Taking
$d = 987$:
\[
\begin{array}{llll}
A: 88.06 \to 88 & B: 5.47 \to 5 & C: 4.38 \to 4 & D: 3.40 \to 3
\end{array}
\]
totaling $100$. \checkmark

Webster's method lands on the same apportionment as Jefferson's here, and so it also
exceeds A's upper quota of $87$. Webster's method is usually the most even-handed of
the three, but on a population this lopsided --- one district holding nearly $87\%$ of
the people --- it can still break the quota rule.

Comparing all four methods on this problem:
\[
\begin{array}{lcccc}
 & A & B & C & D\\
\text{Hamilton} & 87 & 6 & 4 & 3\\
\text{Jefferson} & 88 & 5 & 4 & 3\\
\text{Adams} & 85 & 6 & 5 & 4\\
\text{Webster} & 88 & 5 & 4 & 3
\end{array}
\]
Only Hamilton's method stays within the quota here --- but as Homework 9 shows,
Hamilton's method has troubles of its own.
\end{feedback}
\end{prompt}
\end{problem}

\end{document}
